\begin{theorem}[Heron square-root sibling consumer lattice]
\label{thm:heron-square-root-sibling-consumer-lattice}
Every accepted $\HeronSquareRootUp$ consumer route used by \autoref{ch:concrete-instances-real-reciprocal-namecert}, \autoref{ch:concrete-instances-square-root-continuity-modulus-namecert}, and \autoref{ch:concrete-instances-dyadic-newton-sqrt-namecert} first reads the Heron square-root packet's dyadic approximation, $\RegSeqRatUp$ readback, $\RealUp$ seal, transport, replay, provenance, and local $\NameCert_{\HeronSquareRootUp}$ rows.
\end{theorem}
\begin{proof}
Project the accepted Heron carrier from the chapter spine. The rows $A,S,R,L,W,D$ expose the dyadic radicand, seed, reciprocal ledger, Heron recurrence window, and contraction estimate before any downstream consumer can inspect the regular-rational readback $Q$ or terminal $\RealUp$ seal $E$.

Apply the Heron reciprocal, continuity, and Newton sibling routes already recorded in this chapter. The real reciprocal forward link routes the reciprocal row through \autoref{ch:concrete-instances-real-reciprocal-namecert}; the nonnegative regular-Cauchy reciprocal-window route keeps the same Heron window visible for square-root continuity consumers; and the Newton sibling lattice keeps the dyadic Newton route attached to the displayed approximation rows.

Now read the downstream chapters. \autoref{ch:concrete-instances-real-reciprocal-namecert}, \autoref{ch:concrete-instances-square-root-continuity-modulus-namecert}, and \autoref{ch:concrete-instances-dyadic-newton-sqrt-namecert} therefore receive the Heron packet only after $Q,E,T,C,P,N$ have preserved readback, Real sealing, transport, replay, provenance, and local naming. No consumer route obtains a host square-root value, a hidden reciprocal, quotient-real equality, selected completion, or ambient completeness row.
\end{proof}
